\subsection{环与域的特征}

我们仅当环 A 具有子域时才定义环 A 的特征。此时，设 f 是 Z 到 A 的唯一环同态，并设 n 是 Z 中生成作为 f 的核之理想的唯一正整数（I，第 111 页）；则整数 n 称为 A 的特征。

设 A 是一个已定义特征的环；则 A 不化为 0。由定理 1，A 中存在唯一的素域子域 P；我们称它为 A 的素子域。由定理 1 的证明，有以下两种可能：

\begin{enumerate}
\def\labelenumi{\alph{enumi})}
\item
  A 的特征为 0，P 同构于 Q，
\item
  \(\mathbf{A}\) 的特征是一个素数 \(p\)，\(\mathbf{P}\) 同构于 \(\mathbf{F}_p\) 。
\end{enumerate}

若 A 的特征为零，则存在 Q 到 A 的唯一环同态；其像为 A 的素子域，包含在 A 的中心中。因此，存在唯一的与环结构相容的 A 的 Q-代数结构。当 A 的特征为素数 p 时，把域 Q 替换为域 \(F_{p}\) 即得相应的性质。

命题 1. --- 设 A 是一个不化为 0 的环。

\begin{enumerate}
\def\labelenumi{\alph{enumi})}
\item
  为使 A 具有特征 0，必要且充分的是，映射 \(x \mapsto n \cdot x\)（A 到自身）对每个整数 \(n \neq 0\) 是双射。
\item
  设 \(p\) 为一个素数。为使 \(\mathbf{A}\) 具有特征 \(p\)，必要且充分的是 \(p \cdot x = 0\) 对所有 \(x \in \mathbf{A}\) 成立。
\end{enumerate}

设 \(f\) 是 \(\mathbf{Z}\) 到 \(\mathbf{A}\) 中的唯一同态；我们有 \(n \cdot x = f(n)x\)，对每个整数 \(n\) 及每个 \(x\) （在 \(\mathbf{A}\) 中）。为使 \(\mathbf{A}\) 具有特征 0，必要且充分的是 \(f\) 可扩展为 \(\mathbf{Q}\) 到 \(\mathbf{A}\) 中的一个同态，即 \(f(n)\) 在 \(\mathbf{A}\) 中可逆，对每个 \(n \neq 0\)（I，第 113 页）；这就证明了 \(a\) )。类似地，为使 \(\mathbf{A}\) 具有特征 \(p\)，必要且充分的是 \(f\) 零化 \(p\mathbf{Z}\)，即 \(f(p) = 0\)，或者也即 \(p \cdot x = 0\) 对所有 \(x \in \mathbf{A}\) 成立；这就证明了 \(b\) 。

设 A 为一个不一定交换的域。A 的中心是一个（交换）域；因此 A 的特征和素子域是有定义的。

注记。--- 1) 设 A 与 A' 为两个不化为 0 的环。假设 A 的特征已定义，并且存在一个从 A 到 A' 的同态 u。A 的素子域在 u 下的像是 A' 的一个子域 P'，与 P 同构，因此也是素域。由此可知 A' 的特征已定义，并且等于 A 的特征。若 A 与 A' 的特征为 0（相应地， \(p \neq 0\) ），映射 u 是 Q 上代数的一个同态（相应地， \(F_{p}\) 上代数的同态）。

\begin{enumerate}
\def\labelenumi{\arabic{enumi})}
\setcounter{enumi}{1}
\item
  注 1 表明，若 A 是特征为 0 的环（相应地， \(p \neq 0\) ），则任何以 A 为子环的环 A'，或 A 关于双边理想 \(\mathfrak{a} \neq \mathbf{A}\) 的任何商环，也具有同样的性质。特别地，若 K 是域，则 K 的每个子域和 K 的每个扩域都与 K 具有相同的特征。
\item
  设 A 是域 K 上的一个不化为 0 的代数。由于映射 \(\lambda \mapsto \lambda \cdot 1\)（从 K 到 A）是环同态，注记 1 表明 A 的特征有定义且等于 K 的特征。
\item
  由于域 \(\mathbf{Q}\) 是无限的，每个特征为0的环都是无限的；因此，每个有限域都有非零特征。
\item
  设 A 是一个不化为 0 的环，其加法群是一个无挠 Z-模，并令 \(B = Q \otimes_{Z} A\) 。映射 \(x \mapsto 1 \otimes x\)（从 A 到 B）是单射（II，第 314 页），因此 A 同构于一个特征为 0 的环的子环。
\end{enumerate}

